\chapter*{The question behind the calculations}
\addcontentsline{toc}{chapter}{The question behind the calculations}
Catalan's constant has a simple definition,
\[
G=\sum_{k=0}^{\infty}\frac{(-1)^k}{(2k+1)^2}.
\]
It is easy to locate on the real line. It is much harder to classify arithmetically. An accurate decimal approximation does not decide whether the value is rational. A recurrence that produces astonishingly close rational approximations does not decide that question either, unless the denominators and the nonzero errors satisfy the right bounds.

This reader follows that distinction through the recovered Operation Komorebi work. Its starting point is a small exact arithmetic program with three historical records dated 23 September 2026. The program checks a Wronskian identity, scans the powers of two dividing binomial packet sums, and computes common divisors of cleared coefficient pairs. The mathematical exposition here is newly assembled for this edition. No earlier written Catalan proof manuscript was recovered.

\begin{claimboundary}
This edition does not establish the irrationality of Catalan's constant. It supplies written arguments for specified recurrence and packet identities, checks external inputs in their exact normalization, reproduces finite arithmetic, and states the uniform estimates that would complete this route.
\end{claimboundary}

The strongest statements in the reader have different kinds of support. The Wronskian has an elementary proof valid for every index. The integrality of the scaled first coefficient follows from an authenticated binomial formula and a written coefficient argument. Even filter indices give nonzero error through a positive integral. Those statements should be stated plainly. The large divisibility scan is exact on its finite grid, and its grid should be stated just as plainly. It does not control all later indices.

\section*{How the route would finish}
If $G=a/b$ were rational, every nonzero form $AG-B$ with integers $A,B$ would have absolute value at least $1/b$. Thus nonzero integer forms tending to zero would contradict rationality. The recurrence gives small real errors. Clearing its rational coefficients multiplies those errors by an arithmetic cost. Common divisors can reduce that cost, but a divisor created by an unnecessarily large clearing factor is not itself a new saving. The crucial quantity is the error after the coefficients have been reduced.

The reader first develops the constant and this rationality obstruction. It then follows the recurrence, the Wronskian, the integer packet and the filter. The last chapters measure primitive content and make the remaining asymptotic inequality explicit. A separate visual guide follows the same route with scientific diagrams and bounded data.

\section*{A dated view of the literature}
The edition date matters. A September 2026 preprint by Zhi-Wei Sun claims an irrationality proof. A later primary note by Dane Wachs identifies an incomplete derivation involving the contribution from the prime two. Both are recorded in their checked versions. This reader does not referee every possible repair of that claim, and it does not infer rationality from a criticism of one argument. No irrationality proof is established by the material checked for this publication.

The older recurrence literature also contains a normalization discrepancy worth exposing. A printed integral prefactor conflicts with the initial value. Evaluating the first integral and comparing exact formal coefficients prevents that mismatch from entering the packet calculations. The source contracts record both the printed formula and the normalization actually used here.

\section*{Evidence and reproducibility}
The source bundle contains the three historical records byte for byte, the current LaTeX, the Python figures, fresh run manifests, and separate mathematical, source and layout audits. The historical default output was reproduced byte for byte. An additional implementation compared every pair in the smaller direct grid, rather than comparing only the observed maxima.

An exploratory sign rule for odd filters failed. That failure is retained with an exact counterexample. The restriction to even filter indices has a written nonvanishing argument, so the failed guess need not be disguised or treated as fatal to the entire construction.

Neither repeated computational agreement nor agreement between reviewers is a universal proof. No proof assistant certification is claimed. The purpose of the release is to make the conclusions, dependencies and remaining work concrete enough to inspect and reproduce.

\medskip
\noindent\textbf{Komorebi P.}\\
Operation Komorebi\\
Reader edition 1.0, 9 October 2026
